\section{Discussion}
\label{sec:discussion}
\paragraph{What the evidence supports}
Fusing the two cameras helps on M3FD: against the matched visible-only detector, which is the stronger single camera there, CFFM-Net v2 gains \gainMfdCffmVtwoOverBestMeanMap{}\,$\pm$\,\gainMfdCffmVtwoOverBestStdMap{}\,mAP and the concat baseline \gainMfdConcatMdOverBestMeanMap{}\,$\pm$\,\gainMfdConcatMdOverBestStdMap{}\,mAP, with the caveat that the single-camera rows are one seed. The fuser design adds a smaller step: v2 beats the matched concat baseline on M3FD by \marMfdVtwoVsConcatMdMeanMap{}\,$\pm$\,\marMfdVtwoVsConcatMdStdMap{}\,mAP with the same sign on all three seeds, but the margin is not significant at $n=3$ (paired $p=\sigMfdP{}$; one-sided sign test $\sigMfdSignP{}$), and the small-object margin of \marMfdVtwoVsConcatMdMeanApSmall{}\,$\pm$\,\marMfdVtwoVsConcatMdStdApSmall{}\,pp is not significant either. v2 costs \latLocalRatioMedian{} times the latency, \latLocalRatioParams{} times the parameters and \latLocalRatioGflops{} times the FLOPs of the concat baseline, which passes the limit of at most twice the baseline latency that was fixed in advance. Under sensor faults on M3FD, v2 is ahead of the concat baseline on thermal-side faults by \fltMfdMarginMeanThermalAvgSevThreeMap{}\,$\pm$\,\fltMfdMarginStdThermalAvgSevThreeMap{}\,pp and behind on visible-side faults by \fltMfdMarginMeanVisibleAvgSevThreeMap{}\,$\pm$\,\fltMfdMarginStdVisibleAvgSevThreeMap{}\,pp, a pattern consistent across three seeds but not tested for significance per fault. The video and runtime experiments support three narrower statements. The Temporal State Memory lowers the KAIST miss rate under thermal noise (\memThermalNoiseImageAllMr{}\,\% to \memThermalNoiseMemoryAllMr{}\,\%) and darkness of the visible camera (\memVisibleDarkImageAllMr{}\,\% to \memVisibleDarkMemoryAllMr{}\,\%). The live runtime keeps delay and GPU memory bounded over ten minutes. And its stream monitor flags a cable pull within \livFaultOnUnpluggedDetectS{}\,s.

\paragraph{What it does not support}
The evidence does not show that CFFM-Net is better on LLVIP, where the margin of \marLlvipVtwoVsConcatMdMeanMap{}\,pp has a spread of \marLlvipVtwoVsConcatMdStdMap{}\,pp, or on KAIST stills, where the baseline has the lower miss rate (\kstConcatAllMr{}\,\% against \kstCffmAllMr{}\,\%). It does not show that the gate makes the model more robust than simpler means: a concat baseline trained with fault augmentation retains \fltMfdConcatFtRetainedSevThreePercent{}\,\% of its clean mAP on M3FD against \fltMfdCffmVtwoMeanRetainedSevThreePercent{}\,\% for v2, and v2 loses heavily to noise faults on LLVIP. It does not show that the reliability branch helps: in a preliminary single-seed run v3 is worse than v2 and than v3 without the branch on both datasets. It does not show that the selective scan, as opposed to the other changes made between round one and v2, is the source of the gain, since the round-one controls with a gated convolution of matched cost land within one seed's noise of the scan. It does not show that the memory gain under thermal noise comes from the recurrent state, because the memory was trained on faulty clips and the detector under it was not, and we ran no stateless adapter as a control. It does not show a tracking benefit from the memory, a benefit from feeding the monitor flags to the model (\livFaultMrChange{}\,points of miss rate in a single run), or real-time operation at the camera rate, since the runtime processes \livSoakThirtyFps{} of a nominal \livSoakThirtyNominalFps{} frames per second. And it does not place our models among the published state of the art: on KAIST our miss rates are above \sotaKaistMrMax{}\,\%, the worst published figure we list, and the LLVIP and M3FD figures of Section~\ref{sec:sota} are in different settings.

\paragraph{Limits}
All detectors are nano scale and trained once per seed with three seeds at most; v3 and its controls have one seed, as do the single-camera baselines and every KAIST result. The scan-versus-gated-convolution control was run at round one and has not been repeated at v2 over seeds. The fault benchmark uses synthetic faults and the test split; on LLVIP it covers one seed only, and the fault-trained baseline is scored at one seed on both datasets. The KAIST evaluation uses the original annotations and our own implementation of the miss rate, so it is not comparable with published figures. The live runtime was run on replayed recordings on the same laptop GPU, not on physical cameras, and its fault script covers three events in one run. Latency in Table~\ref{tab:latency} was measured on one \latLocalDevice{}, so absolute values should not be compared with other papers. The multi-camera layer ran on RGB WILDTRACK, so it supports the claim that the system's last stage works, and makes no claim about fusion.

\paragraph{Next steps}
Four experiments follow from the results. First, run the fault benchmark on a scan-based model trained with fault augmentation against the fault-trained concat baseline, at three seeds on both datasets, which decides whether the gate adds robustness once training is matched. Second, train a stateless adapter of the memory's size on the same clips, to separate the recurrent state from the extra capacity. Third, repeat the scan-versus-convolution comparison in the v2 recipe at three seeds. Fourth, find the cause of the loss to additive noise on LLVIP, whose gate has no cue that separates noise from texture.
